% ===== PRÉAMBULE CANONIQUE — à coller VERBATIM en tête de CHAQUE .tex =====
\documentclass[11pt,a4paper]{article}
\usepackage[utf8]{inputenc}\usepackage[T1]{fontenc}\usepackage{lmodern}
\usepackage[french]{babel}
\usepackage{amsmath,amssymb,mathtools}\usepackage{icomma}
\usepackage[version=4]{mhchem}          % \ce{...} : équations & formules
\usepackage{siunitx}\sisetup{locale=FR} % \qty{}{}, \si{}, \ang{}
\usepackage{chemfig}                    % \chemfig{...} : structures développées
\usepackage{xcolor}\usepackage{colortbl}
\usepackage{tikz}\usepackage{pgfplots}\pgfplotsset{compat=1.18}
\usepackage[most]{tcolorbox}\usepackage{enumitem}\usepackage{array,booktabs}
\usepackage[a4paper,margin=2.2cm]{geometry}
\usetikzlibrary{arrows.meta,calc,angles,quotes,positioning,patterns,backgrounds,shapes.geometric}
\definecolor{primary}{RGB}{222,49,99}\definecolor{primarydark}{RGB}{150,22,55}
\definecolor{defcol}{RGB}{0,114,178}\definecolor{methcol}{RGB}{52,92,148}
\definecolor{excol}{RGB}{0,135,90}\definecolor{attncol}{RGB}{200,40,55}
\tcbset{boxbase/.style={enhanced,breakable,boxrule=0.9pt,arc=2.5pt,
  left=3mm,right=3mm,top=2.3mm,bottom=2.3mm,fonttitle=\bfseries,coltitle=white}}
% --- boîtes TUTORIEL (noms EXACTS, ne pas renommer) ---
\newtcolorbox{cle}[1][]{boxbase,colback=defcol!6,colframe=defcol,
  title={\textsc{À retenir}\ifx\relax#1\relax\relax\else\textmd{ : \itshape #1}\fi}}
\newtcolorbox{methode}[1][]{boxbase,colback=methcol!7,colframe=methcol,sharp corners,
  title={\textsc{Méthode}\ifx\relax#1\relax\relax\else\textmd{ --- \itshape #1}\fi}}
\newtcolorbox{exemplebox}[1][]{boxbase,colback=excol!5,colframe=excol,
  title={\textsc{Exemple}\ifx\relax#1\relax\relax\else\textmd{ : \itshape #1}\fi}}
\newtcolorbox{piege}[1][]{boxbase,colback=attncol!6,colframe=attncol,
  title={\textsc{Piège}\ifx\relax#1\relax\relax\else\textmd{ : \itshape #1}\fi}}
\newtcolorbox{remarque}[1][]{boxbase,colback=black!4,colframe=black!45,
  title={\textsc{Remarque}\ifx\relax#1\relax\relax\else\textmd{ : \itshape #1}\fi}}
% --- boîtes EXERCICE / CORRIGÉ (noms EXACTS, auto-comptées) ---
\newtcolorbox[auto counter]{exo}[1][]{boxbase,colback=primary!4,colframe=primary,
  title={\textsc{Exercice}~\thetcbcounter\ifx\relax#1\relax\relax\else\textmd{ --- \itshape #1}\fi}}
\newtcolorbox[auto counter]{corrige}[1][]{boxbase,colback=excol!4,colframe=excol,
  title={\textsc{Corrigé}~\thetcbcounter\ifx\relax#1\relax\relax\else\textmd{ --- \itshape #1}\fi}}
\newcommand{\R}{\mathbb{R}}\newcommand{\N}{\mathbb{N}}\newcommand{\Z}{\mathbb{Z}}\newcommand{\ds}{\displaystyle}
% (un \usepackage{fancyhdr}+en-tête est possible ; il est ignoré côté web)
% =========================================================================
\begin{document}

\begin{center}{\large\bfseries\color{primarydark} Thème 4 — Corrigé Facile}\end{center}

\begin{corrige}[écrire la dissociation]
\begin{enumerate}[label=\alph*)]
  \item \ce{NaCl -> Na+ + Cl-} \ ($1$ \ce{Na+}, $1$ \ce{Cl-}).
  \item \ce{CaCl2 -> Ca^2+ + 2 Cl-} \ ($1$ \ce{Ca^2+}, $2$ \ce{Cl-}).
  \item \ce{Na2SO4 -> 2 Na+ + SO4^{2-}} \ ($2$ \ce{Na+}, $1$ \ce{SO4^{2-}}).
  \item \ce{Al2(SO4)3 -> 2 Al^3+ + 3 SO4^{2-}}.
  \item \ce{K3PO4 -> 3 K+ + PO4^{3-}}.
\end{enumerate}
\end{corrige}

\begin{corrige}[concentration des ions]
\begin{enumerate}[label=\alph*)]
  \item $[\ce{Na+}] = \SI{0.5}{\mole\per\litre}$, $[\ce{Cl-}] = \SI{0.5}{\mole\per\litre}$.
  \item $[\ce{Ca^2+}] = \SI{0.2}{\mole\per\litre}$, $[\ce{Cl-}] = 2\times 0{,}2 = \SI{0.4}{\mole\per\litre}$.
  \item $[\ce{Na+}] = 2\times 0{,}1 = \SI{0.2}{\mole\per\litre}$, $[\ce{SO4^{2-}}] = \SI{0.1}{\mole\per\litre}$.
\end{enumerate}
\end{corrige}

\begin{corrige}[quantité de matière d'un ion]
\begin{enumerate}[label=\alph*)]
  \item \ce{CaCl2} libère $2$ \ce{Cl-} : $n(\ce{Cl-}) = 2\times 0{,}3\times 1 = \SI{0.6}{\mole}$.
  \item \ce{Na3PO4} libère $3$ \ce{Na+} : $n(\ce{Na+}) = 3\times 0{,}1\times 2 = \SI{0.6}{\mole}$.
\end{enumerate}
\end{corrige}

\begin{corrige}[quantité ou concentration~?]
\begin{enumerate}[label=\alph*)]
  \item C'est une \emph{quantité} : $n(\ce{NO3-}) = 2\times 0{,}05 = \SI{0.1}{\mole}$, à diviser par
        $V = \SI{0.2}{\litre}$ : $[\ce{NO3-}] = \dfrac{0{,}1}{0{,}2} = \SI{0.5}{\mole\per\litre}$.
  \item C'est déjà une \emph{concentration} : $[\ce{NO3-}] = 2\times 0{,}05 = \SI{0.1}{\mole\per\litre}$
        (\textbf{sans} diviser par le volume).
\end{enumerate}
\end{corrige}

\begin{corrige}[la plus riche en ions chlorure]
$[\ce{Cl-}] = \text{(nombre de \ce{Cl})}\times 0{,}3$ :
\begin{enumerate}[label=\alph*)]
  \item \ce{NaCl} : $[\ce{Cl-}] = \SI{0.3}{\mole\per\litre}$.
  \item \ce{CaCl2} : $[\ce{Cl-}] = 2\times 0{,}3 = \SI{0.6}{\mole\per\litre}$.
  \item \ce{AlCl3} : $[\ce{Cl-}] = 3\times 0{,}3 = \SI{0.9}{\mole\per\litre}$.
\end{enumerate}
Ordre croissant : \ce{NaCl} $<$ \ce{CaCl2} $<$ \ce{AlCl3}.
\end{corrige}

\end{document}
